\documentclass[11pt]{article}
\usepackage{fontspec}
\setmainfont{Times New Roman}
\setsansfont{Arial}
\setmonofont{Consolas}
\usepackage{amsmath,amssymb}
\usepackage{geometry}
\usepackage{microtype}
\geometry{a4paper,margin=3cm}
\setlength{\parindent}{1.25em}
\setlength{\parskip}{0pt}
\emergencystretch=10em
\allowdisplaybreaks
\providecommand{\Kfield}{\mathfrak{K}}
\providecommand{\Ssys}{\mathfrak{S}}
\providecommand{\Mfield}{\mathfrak{M}}
\providecommand{\tq}{\tau}
\providecommand{\ds}{\displaystyle}

\begin{document}

\begin{center}
{\Large\bfseries 5. Поля рациональных функций}\par
\vspace{1.2em}
J. Ber. d. DMV 22 (1913), S. 316--319
\end{center}

\vspace{2em}

Следующие вопросы первоначально восходят к беседам с Э. Фишером. Некоторые из них, впрочем, для специального случая одной неопределенной -- причем при более общих предположениях о поле коэффициентов -- уже были поставлены и решены Э. Штайницем.\footnote{E. Steinitz: Algebraische Theorie der Körper (особенно \S\ 24). Crelles Journal 137. 1910.}

Под ``полем рациональных функций'' я понимаю поле, элементы которого являются рациональными функциями от $n$ неопределенных с коэффициентами из заранее заданного числового поля, которое, в частности, может содержать и все комплексные числа. Примерами таких полей служат поле симметрических функций от $n$ величин или, более общо, совокупность рациональных функций от $n$ неопределенных, допускающих перестановки некоторой группы (лагранжевы области родов); далее -- поле инвариантов.

Между двумя первыми названными полями и последним имеется характерное различие: первые содержат $n$ алгебраически независимых функций, элементарные симметрические функции; тогда как число алгебраически независимых инвариантов всегда меньше числа неопределенных. Поэтому важно, что вообще такие поля второго типа можно взаимно однозначно сопоставлять с полями первого типа, заменяя некоторые из неопределенных числами. Поэтому далее я ограничусь полями первого типа, то есть полями с $n$ алгебраически независимыми функциями; в силу этого сопоставления результаты затем имеют общий характер.

Вопросы группируются вокруг трех базисных понятий: рациональная база, минимальная база и интегральная база.

1. Под ``рациональной базой'' я понимаю конечное число функций поля такое, что всякую функцию поля можно представить как рациональную комбинацию этого конечного числа функций с коэффициентами из заданной области коэффициентов. Простыми рассуждениями -- которые в случае одной неопределенной имеются также у Э. Штайница, там же -- доказывается существование рациональной базы для всякого поля рациональных функций. Например, по теореме Лагранжа элементарные симметрические функции и одна функция, принадлежащая группе, образуют рациональную базу для упомянутых выше лагранжевых областей родов.

2. Всякая рациональная база должна (для полей первого типа) содержать по меньшей мере $n$ функций; если она содержит ровно $n$ функций, которые тогда должны быть алгебраически независимыми, я называю ее ``минимальной базой''. Вопрос о существовании минимальной базы частично решается теоремами Люрота, Кастельнуово и Энрикеса.\footnote{J. Lüroth: Beweis eines Satzes über rationale Kurven. Math. Ann. 9. 1875. --- G. Castelnuovo: Sulla razionalità delle involuzioni piane. Math. Ann. 44. 1893. --- F. Enriques: Sopra una involuzione non razionale dello spazio. Rend. Acc. Linc. vol. 21. 21. Jan. 1912.} Из них следует, что для полей от одной неопределенной минимальная база всегда существует: это функция Люрота данного поля, единственно определенная с точностью до дробно-линейного преобразования (ср. Штайниц, \S\ 24). Для полей от двух неопределенных минимальная база существует всегда и, вообще говоря, только тогда, когда допускается алгебраическое расширение поля коэффициентов, тогда как для трех и более неопределенных минимальная база вообще не существует. Охарактеризовать специальные поля с минимальной базой еще не пытались.

Теперь я продолжила исследовать вопрос о минимальной базе для лагранжевых областей родов. Здесь он приобретает следующий смысл: ``Если лагранжева область родов, принадлежащая группе $G$, обладает минимальной базой, то всю совокупность уравнений с заданной группой $G$ можно рационально построить посредством параметрического представления; и именно для всякого числового поля, содержащего область коэффициентов минимальной базы -- следовательно, в частности, для любого числового поля, если эта область коэффициентов есть поле рациональных чисел.''

Самый простой пример дает симметрическая группа; здесь элементарные симметрические функции образуют минимальную базу с рациональными числовыми коэффициентами. Отсюда в качестве параметрического представления получается просто уравнение с неопределенными коэффициентами. Как от него перейти к сколь угодно многим уравнениям ``без аффекта'' над каждым числовым полем, показывает теорема Гильберта о неприводимости.

Теперь непосредственно из теорем Люрота и Кастельнуово следует существование минимальной базы для всех групп, возникающих при уравнениях 3-й и 4-й степеней; при этом далее оказывается, что эта минимальная база имеет рациональные числовые коэффициенты. Значит, над произвольным числовым полем можно рационально построить совокупность уравнений 3-й и 4-й степеней с заданной группой. Для циклических и диэдральных групп существует минимальная база с полем $n$-х корней из единицы в качестве области коэффициентов; то же справедливо и для некоторых других метациклических групп. Вопрос, является ли минимальная база вообще характерной для некоторых групп, должен остаться нерешенным.

3. Чтобы перейти к последнему базисному понятию, я рассматриваю совокупность полиномов, содержащихся в поле. Под ``интегральной базой'' я понимаю конечное число этих полиномов такое, что всякий полином поля можно представить как целую рациональную комбинацию этого конечного числа полиномов с коэффициентами из заданной области коэффициентов. Вопрос о существовании интегральной базы -- который, например, как специальный случай включает конечность системы инвариантов -- был поставлен Гильбертом в его ``Математических проблемах''\footnote{D. Hilbert: Mathematische Probleme. Доклад, Париж 1900. Проблема 14. Göttinger Nachr. 1900.} в несколько иной формулировке, как проблема относительно целых функций.

Пока я могу указать лишь один класс полей, для которых интегральная база существует. А именно, если поле содержит систему из $n$ полиномов такую, что однородная результанта членов наивысшей степени этих полиномов отлична от нуля, то интегральная база существует; ее дают коэффициенты всех $u$ в неприводимом над полем уравнении для линейной формы:
\[
  u_0=u_1x_1+u_2x_2+\cdots+u_nx_n.\footnotemark
\]
\footnotetext{Это неприводимое уравнение в случае одной неопределенной встречается у Э. Штайница в доказательстве теоремы Люрота.}
Если, в частности, значение результанты равно единице области коэффициентов, то обеспечена также интегральность представления. Пример этого класса полей дают лагранжевы области родов, поскольку результанта элементарных симметрических функций имеет значение $\pm 1$. Еще один пример (без интегральности) дают все поля, содержащие только один алгебраически независимый полином; интегральная база здесь тождественна функции Люрота наименьшего подполя, содержащего все полиномы. Так, как известно, все целые рациональные проективные инварианты квадратичной формы от $n$ переменных являются целыми функциями дискриминанта.

\enlargethispage{2\baselineskip}
Указанное условие является только достаточным, отнюдь не необходимым для существования интегральной базы. Легко построить поля, в которых результанта любых $n$ полиномов обращается в нуль и которые тем не менее обладают интегральной базой, заданной коэффициентами того неприводимого уравнения. Возникает вопрос, дают ли, быть может, и в самом общем случае коэффициенты того уравнения интегральную базу.

\end{document}
